% This file compiles with both LuaLaTeX and XeLaTeX
\documentclass[11pt]{article}

% Change "review" to "final" to generate the final (sometimes called camera-ready) version.
% Change to "preprint" to generate a non-anonymous version with page numbers.
\usepackage[review]{acl}
\usepackage{algorithm}
\usepackage{algorithmic}
% This is not strictly necessary, and may be commented out,
% but it will improve the layout of the manuscript,
% and will typically save some space.
 \usepackage{microtype}
 \usepackage{amsmath}

% If the title and author information does not fit in the area allocated, uncomment the following
%
%\setlength\titlebox{<dim>}
%
% and set <dim> to something 5cm or larger.

% These font selection commands work with
% LuaLaTeX and XeLaTeX, but not pdfLaTeX.
\usepackage[english,bidi=default]{babel} % English as the main language.
\babelfont{rm}{TeXGyreTermesX} % similar to Times
%%% include whatever languages you need below this line
\babelprovide[import]{hindi}
\babelfont[*devanagari]{rm}{Lohit Devanagari}
\babelprovide[import]{arabic}
\babelfont[*arabic]{rm}{Noto Sans Arabic}


%\usepackage{polyglossia}
%\setdefaultlanguage{english}
%\setotherlanguages{arabic,russian,thai,hindi,kannada}

%%%%%


\title{LuaLaTeX and XeLaTeX Template for *ACL Style Files}

% Author information can be set in various styles:
% For several authors from the same institution:
% \author{Author 1 \and ... \and Author n \\
%         Address line \\ ... \\ Address line}
% if the names do not fit well on one line use
%         Author 1 \\ {\bf Author 2} \\ ... \\ {\bf Author n} \\
% For authors from different institutions:
% \author{Author 1 \\ Address line \\  ... \\ Address line
%         \And  ... \And
%         Author n \\ Address line \\ ... \\ Address line}
% To start a seperate ``row'' of authors use \AND, as in
% \author{Author 1 \\ Address line \\  ... \\ Address line
%         \AND
%         Author 2 \\ Address line \\ ... \\ Address line \And
%         Author 3 \\ Address line \\ ... \\ Address line}

\author{First Author \\
  Affiliation / Address line 1 \\
  Affiliation / Address line 2 \\
  Affiliation / Address line 3 \\
  \texttt{email@domain} \\\And
  Second Author \\
  Affiliation / Address line 1 \\
  Affiliation / Address line 2 \\
  Affiliation / Address line 3 \\
  \texttt{email@domain} \\}

\begin{document}

\maketitle
\begin{abstract}
This document provides an example showing how
to use the *ACL style files with either
LuaLaTeX or XeLaTeX.
\end{abstract}


\section{Introduction}

Please see the general instructions
in the file \verb|acl_latex.tex|.

Here are some examples of text in various languages.

Hindi: \foreignlanguage{hindi}{मानव अधिकारों की सार्वभौम घोषणा}

Arabic: \foreignlanguage{arabic}{الإعلان العالمي لحقوق الإنسان}

Here is an example citation:
\citet{Gusfield:97} argues that...


\section{Related Work}
\label{sec:related}

\noindent \textbf{Transformer-based fusion}.
MulT \citep{tsai2019multimodal} stacks cross-modal attention blocks for unaligned
MSA; MAG-BERT \citep{rahman2020integrating} injects nonverbal context into
language models.
These models typically form a \emph{single} fused representation before the head.
InfoGate instead keeps an explicit primary vs.\ auxiliary asymmetry after IB
compression.

\noindent \textbf{Modality-specific representation learning}.
MISA \citep{hazarika2020misa} disentangles invariant and specific factors; Self-MM
\citep{yu2021learning} adds self-supervision.
InfoGate is complementary: we do not maintain parallel invariant/specific
towers, but we \emph{do} enforce per-modality bottlenecks with shared
dimensionality before fusion.

\noindent \textbf{Information bottlenecks and debiasing}.
Deep variational IB \citep{alemi2017deep} motivates stochastic encoders; CaMIB
\citep{jiang2025towards} combines IB-style filtering with causal masks for OOD
robustness.
InfoGate is \emph{not} a causal identification paper; rather, we use IB
confidence as a \emph{mechanistic} control signal inside attention, closer to
uncertainty-aware fusion than to full causal graphs.

\noindent \textbf{Dynamic primary modality selection}.
Lines of work on modality optimization and dynamic primaries
\citep{mods2026align,yang2025improving} argue that the ``best'' carrier modality varies by sample.
Our MSelector realizes this with lightweight pooling scores and Gumbel--Softmax
gradients, with stage-two KL targets derived from per-modality label prediction
errors (``select the modality that already fits the label'').

\noindent \textbf{Conceptual connections to hierarchical and token-level IB}.
Hierarchical IB fusion \citep{ithp2024align} and token-level IB perspectives
\citep{cyin2025align} informed exploratory designs; we cite them as internal
alignment references and emphasize that InfoGate's implementation is
self-contained and should be evaluated on its own merits.

\section{Methodology}
\label{sec:methodology}

\subsection{Positioning relative to MODS and CyIN}
\label{sec:mods-cyin-positioning}

\textbf{MODS.}
MODS \citep{mods2026align} advocates modality optimization with a
\emph{sample-dependent} primary stream: the model should not assume that text
always carries the decisive evidence, and fusion should be organized around
whichever modality best supports the label for each instance.
InfoGate adopts the same high-level stance through MSelector
(\S\ref{sec:method-msel}), but instantiates it on \emph{IB-compressed}
bottlenecks rather than on raw modality streams, and routes the final
prediction exclusively through the enhanced primary bottleneck (no direct
textual shortcut into the head), which matches the primary-centric design
discipline emphasized in recent modality-optimization work
\citep{yang2025improving}.

\textbf{CyIN.}
CyIN \citep{cyin2025align} motivates token-level information bottlenecks and
informative latent spaces that suppress nuisance variation while retaining
task-relevant structure.
InfoGate is aligned with that goal in that each modality passes through a
stochastic bottleneck whose variance supplies a per-token confidence map, but
differs in \emph{where} that signal is consumed: we use IB uncertainty to bias
cross-modal attention and gates inside InfoGate layers (\S\ref{sec:method-layer}),
rather than treating the bottleneck solely as a standalone encoder objective.

\textbf{Synthesis.}
Together, MODS-style primary routing and CyIN-style token-level compression
inform InfoGate's training story---IB warm-up, differentiable primary selection,
routing supervision, and confidence-aware fusion---but the implementation and
evaluation protocol below are self-contained.

\subsection{Problem Statement}
\label{sec:method-prob}
The multimodal sentiment analysis (MSA) task predicts a continuous or discrete
sentiment target from three modalities---language ($t$), acoustic ($a$), and
visual ($v$)---extracted from short video clips.
Given a training set $\{(X^{(i)}_t,X^{(i)}_a,X^{(i)}_v,y^{(i)})\}_{i=1}^{N}$
with labels $y^{(i)}\in[-3,3]$ on CMU-MOSI and CMU-MOSEI and dataset-specific
ranges on CH-SIMS v2, we map every modality to a common hidden size $d$ and
denote the unified unimodal sequence features as
$F_m \in \mathbb{R}^{B \times T \times d}$ with token mask
$M_m \in \{0,1\}^{B \times T}$ for $m \in \{t, a, v\}$.
Our goal is to learn a predictor
$f_{\theta}:\ (X_t, X_a, X_v) \mapsto \hat{y}$ with parameters $\theta$ by
minimizing the empirical risk $\tfrac{1}{N}\sum_i \mathcal{L}^{(i)}$, where
$\mathcal{L}^{(i)}$ is specified in \S\ref{sec:method-loss}.
At inference, routing uses the argmax over selector logits (no Gumbel noise).

\subsection{Framework Overview}
\label{sec:method-overview}
The overall architecture of \textbf{InfoGate} is illustrated in
Figure~\ref{fig:arch} and summarized in Algorithm~\ref{alg:forward}.
First, each modality is passed through a lightweight projection followed by
an IB encoder that outputs a bottleneck representation $B_m$ together with a
per-token confidence signal $C_m$ derived from the bottleneck variance.
The three bottleneck sequences are then fed into a Primary Modality Selector
(\textbf{MSelector}) to determine a per-sample primary modality $p$ and two
secondary modalities $a_1,a_2$, together with their routed bottlenecks and
confidences $(B_p, C_p)$, $(B_{a_1}, C_{a_1})$ and $(B_{a_2}, C_{a_2})$.
After primary modality selection, the three streams are progressively processed
by a stack of \textbf{InfoGate layers} that carry out confidence-biased
cross-modal attention and adaptive gating centered on the primary modality,
yielding an enhanced primary bottleneck $B_p^{\mathrm{out}}$.
Finally, $B_p^{\mathrm{out}}$ is aggregated by a masked pooling operation and
passed through a multilayer perceptron (MLP) head to obtain the sentiment
prediction $\hat{y}$.



\begin{algorithm}[t]
  \caption{InfoGate forward pass on a single batch.}
  \label{alg:forward}
  \begin{algorithmic}[1]
    \STATE Project inputs to $F_t,F_a,F_v$ and encode IB latents
      $(B_m,\mu_m,\log v_m,C_m)$ for $m\in\{t,a,v\}$.
    \STATE Compute per-modality IB objective
      $\mathcal{L}_{\mathrm{vib}}$ and pooled label-supervised term
      $\mathcal{L}_{\mathrm{pib}}$.
    \STATE Obtain routing weights and hard choice
      $(w,y)\leftarrow \mathrm{MSelector}(B_t,B_a,B_v)$ via Gumbel--Softmax.
    \STATE Route $(B_p,C_p,B_{a_1},C_{a_1},B_{a_2},C_{a_2})$ from
      $(B_m,C_m,w,y)$ with straight-through gradients.
    \STATE $B_p^{\mathrm{out}} \leftarrow
      \mathrm{InfoGateStack}(B_p,C_p,B_{a_1},C_{a_1},B_{a_2},C_{a_2})$.
    \STATE $\hat{y} \leftarrow \mathrm{Head}(\mathrm{Pool}(B_p^{\mathrm{out}}))$.
    \IF{stage $= 2$}
      \STATE Compute routing supervision
        $\mathcal{L}_{\mathrm{route}}$ from per-modality prediction errors.
    \ENDIF
    \RETURN $\hat{y}$,
      $\mathcal{L}_{\mathrm{vib}}+\mathcal{L}_{\mathrm{pib}}$,
      $\mathcal{L}_{\mathrm{route}}$.
  \end{algorithmic}
\end{algorithm}

\subsection{Unimodal Projections and Variational IB Encoders}
\label{sec:method-ibenc}
Learned linear layers first project text, acoustic, and visual inputs to the
unified feature tensors $F_m \in \mathbb{R}^{B \times T \times d}$.
Each $F_m$ is then passed through an IB encoder, realized as a small MLP whose
output head predicts a Gaussian mean and log-variance
$\mu_m, \log v_m \in \mathbb{R}^{B \times T \times d_b}$, where $d_b$ is the
bottleneck width. Using the reparameterization trick, we draw
\begin{equation}
  B_m = \mu_m + \epsilon \odot \exp\!\bigl(\tfrac{1}{2}\log v_m\bigr),
  \qquad
  \epsilon \sim \mathcal{N}(0, I),
\end{equation}
during training and set $B_m = \mu_m$ at inference.
From the same variance, we derive a per-token confidence map
\begin{equation}
  C_m = \sigma\!\bigl(-\log v_m\bigr) \in (0, 1)^{B \times T \times d_b},
\end{equation}
so that high-variance positions receive low confidence and are naturally
attenuated in subsequent fusion.

\paragraph{Variational IB objective.}
To encourage each bottleneck to discard modality-specific nuisance while
retaining information needed to reconstruct the upstream feature, we introduce
modality-specific decoders $g_m$ that map bottlenecks back to the unified
space and minimize
\begin{equation}
  \mathcal{L}_{\mathrm{vib}}
  = \sum_{m}
    \mathbb{E}\!\left[
      \mathrm{KL}\!\bigl(q(B_m\mid F_m)\,\|\,\mathcal{N}(0,I)\bigr)
      + \beta_{\mathrm{IB}}\,\bigl\|g_m(B_m) - F_m\bigr\|_2^2
    \right],
\end{equation}
where the KL and reconstruction terms are averaged over valid tokens via the
modality mask $M_m$, and $\beta_{\mathrm{IB}}$ trades compression against
fidelity.

\paragraph{Pooled label-supervised IB prediction.}
To inject high-level semantic supervision into each bottleneck, we further
attach per-modality linear probes $\psi_m$ on the masked-mean pooled
bottleneck $\bar{B}_m$. Let $e_m = |\psi_m(\bar{B}_m) - y|$ denote the
per-modality prediction error (or BCE loss in the classification case).
We aggregate the probe losses and pooled-KL terms into
\begin{equation}
  \mathcal{L}_{\mathrm{pib}}
  = \frac{1}{|\mathcal{M}|}\sum_{m \in \mathcal{M}}
    \Bigl(e_m + \beta_{\mathrm{IB}}\,\mathrm{KL}\!\bigl(q(\bar{B}_m)\,\|\,\mathcal{N}(0,I)\bigr)\Bigr),
\end{equation}
with $\mathcal{M} = \{t, a, v\}$. The per-modality errors $e_m$ are also reused
as routing targets in \S\ref{sec:method-rib}.

\subsection{Dynamic Primary Modality Selection}
\label{sec:method-msel}
The sample-wise primary modality is determined by an adaptive MSelector
module, similar in role to the dynamic selector in MODS
\citep{mods2026align} but operating on IB bottleneck sequences rather than raw
features.
For each modality, a scalar projection followed by softmax attention over
time aggregates $B_m$ into a one-dimensional vector
\begin{equation}
\begin{split}
a_m &= \mathrm{softmax}(B_m W_m/\sqrt{d_b})^{\top},\\
h_m &= a_m B_m,\quad m\in\{t,a,v\}.
\end{split}
\end{equation}
A shallow MLP then maps $\mathrm{concat}(h_t, h_a, h_v)$ to routing logits
$z \in \mathbb{R}^3$. During training we use Gumbel--Softmax with the
straight-through trick \citep{jang2017gumbel} to obtain a hard one-hot primary
choice together with a soft routing weight:
\begin{equation}
  y = \mathrm{GumbelSoftmax}\!\bigl(z,\tau_{\mathrm{gumbel}},\mathrm{hard}\bigr),
  \qquad
  w = \mathrm{softmax}\!\bigl(z / \tau_{\mathrm{gumbel}}\bigr),
\end{equation}
with temperature $\tau_{\mathrm{gumbel}}$ annealed linearly across epochs.
The primary bottleneck is assembled as a straight-through mixture
$B_p = \sum_m y_m B_m$, while the two secondary bottlenecks
$B_{a_1}, B_{a_2}$ are the remaining modalities weighted by their soft routing
scores. Their associated confidence maps $C_p, C_{a_1}, C_{a_2}$ are routed
identically, so downstream attention sees confidence that is consistent with
the selected primary.

\subsection{Information-Bottleneck-Guided Fusion}
\label{sec:method-layer}
Stacked InfoGate layers refine the primary stream using the routed secondary
streams, playing a role analogous to the primary-modality-centric fusion
blocks in MODS but replacing plain cross-attention with IB-guided attention
\citep{alemi2017deep}.

\paragraph{IB-guided attention.}
For a multi-head attention block, let $Q$, $K$, $V$ be the standard linear
projections, so that the raw scaled-dot-product scores are
$S = (Q K^{\top})/\sqrt{d/h}$. We inject the per-token confidence of the key
sequence as an additive log-confidence bias on the attention logits and as a
multiplicative gate on the values:
\begin{equation}
  \tilde{S}_{i,j}
    = S_{i,j} + \gamma \cdot \log\!\bigl(\bar{c}_j + \epsilon\bigr),
  \qquad
  \tilde{V}_j = V_j \odot \bar{c}_j,
\end{equation}
where $\gamma$ is a learnable scalar and $\bar{c}_j$ is the confidence of
position $j$ averaged over the bottleneck channel dimension.
High-confidence keys therefore receive larger logits and contribute values
closer to their original magnitudes, while low-confidence (high-variance)
keys are softly suppressed in both the attention distribution and the value
aggregation.

\paragraph{Layer mechanics.}
Each InfoGate layer performs three operations in parallel on the primary
bottleneck: (i) IB-guided cross-attention from primary queries to each of the
two secondary modalities, with their confidence as above; (ii) IB-guided
self-attention within the primary stream; and (iii) cosine-similarity
alignment factors that modulate an adaptive gate
$g = \sigma\bigl(W\,[\mathrm{LN}(B_p)\,\|\,\mathrm{CA}(B_p)]\bigr)$ on the
cross-attention outputs. The three contributions are summed with a residual
connection and followed by a position-wise feed-forward block. All but the
last layer also update the two secondary streams via a symmetric reverse
cross-attention so that auxiliary tokens absorb primary context; the final
layer retains only the primary output $B_p^{\mathrm{out}}$, which is the
representation passed to the prediction head.

\subsection{Routing Supervision}
\label{sec:method-rib}
Stage~1 of training lets the selector drift freely so that the IB encoders can
warm up. In stage~2 we add a routing supervision term that aligns the soft
routing weights $w$ with per-modality prediction errors from the label-level
probes. Specifically, for each modality we convert the error $e_m$ into a
soft target
\begin{equation}
  t_m \propto \exp\!\bigl(-e_m / \tau_{\mathrm{sel}}\bigr),
  \qquad
  \sum_{m} t_m = 1,
\end{equation}
and define the routing loss as a masked KL divergence
\begin{equation}
  \mathcal{L}_{\mathrm{route}}
    = \mathbb{I}[\,\Delta_e \ge \tau_{\mathrm{gap}}\,]\cdot
      \mathrm{KL}\bigl(w\,\|\,t\bigr)
    + \lambda_{\mathrm{ent}}\,\mathcal{H}(w),
\end{equation}
where $\Delta_e$ is the spread of per-modality errors in the batch and
$\tau_{\mathrm{gap}}$ is a gap threshold that prevents the KL term from
firing when all modalities already fit the label comparably. The optional
entropy regularizer $\mathcal{H}(w)$, weighted by $\lambda_{\mathrm{ent}}$,
discourages the selector from collapsing to a single static modality.

\subsection{Training Objective Optimization}
\label{sec:method-loss}
The sentiment prediction is obtained as
$\hat{y} = \mathrm{Head}(\mathrm{LN}(\mathrm{Pool}(B_p^{\mathrm{out}})))$
with a masked-mean pooling operator and a linear regression head.
We adopt a hybrid task loss that combines mean absolute error with a
mean-squared-error regularizer,
\begin{equation}
  \mathcal{L}_{\mathrm{task}}
    = |\hat{y} - y| + \rho_{\mathrm{mse}}\,(\hat{y} - y)^2,
\end{equation}
where $\rho_{\mathrm{mse}}$ weights the MSE component.
The overall training objective unifies the task loss, the variational and
pooled IB objectives, and the stage-two routing supervision:
\begin{equation}
  \mathcal{L}
    = \mathcal{L}_{\mathrm{task}}
    + \alpha_{\mathrm{IB}}\!\left(\mathcal{L}_{\mathrm{vib}}
        + \mathcal{L}_{\mathrm{pib}}\right)
    + \lambda_{\mathrm{route}}\,\mathcal{L}_{\mathrm{route}},
\end{equation}
with $\mathcal{L}_{\mathrm{route}} \equiv 0$ in stage~1 and activated in
stage~2.
Compared with symmetric MulT-style stacks \citep{tsai2019multimodal}, the
design of InfoGate \emph{asymmetrizes} the fusion: only the primary stream
carries the final prediction, and keys with high IB uncertainty are
down-weighted, while the per-layer cost stays in the same
$\mathcal{O}(L \cdot T^2 d)$ regime as a standard Transformer block.




\section{Experiments}

% Entries for the entire Anthology, followed by custom entries
\bibliography{custom}

\appendix

\section{Example Appendix}
\label{sec:appendix}

This is an appendix.

\end{document}
